\documentclass[12pt,a4paper]{report}

\usepackage{fontspec}
\setmainfont{Times New Roman}
\usepackage{geometry}
\geometry{left=3cm,right=2cm,top=2.5cm,bottom=2.5cm}
\usepackage{setspace}
\onehalfspacing
\usepackage{booktabs}
\usepackage{longtable}
\usepackage{array}
\usepackage{hyperref}
\usepackage{xcolor}
\usepackage{listings}
\usepackage{amsmath}
\usepackage{amssymb}

\hypersetup{
    colorlinks=true,
    linkcolor=blue,
    urlcolor=blue,
    citecolor=blue
}

\lstset{
    basicstyle=\ttfamily\small,
    breaklines=true,
    frame=single,
    columns=fullflexible
}

\begin{document}

\begin{titlepage}
    \centering
    {\Large \textbf{BÁO CÁO DỰ ÁN BIG DATA}}\\[1.5cm]
    {\Huge \textbf{Hệ Thống Phát Hiện Thao Túng Thị Trường}}\\[0.4cm]
    {\Large \textbf{Market Manipulation Detection System}}\\[1.2cm]
    {\large Kết hợp C++ Limit Order Book, Machine Learning, Apache Spark và Shapley Value}\\[2cm]
    \vfill
    {\large Ngày hoàn thiện: 2026}\\[0.5cm]
\end{titlepage}

\tableofcontents
\newpage

\chapter{Tổng Quan Dự Án}

\section{Bối cảnh}

Thị trường tài chính hiện đại, đặc biệt là thị trường tiền điện tử và các hệ thống giao dịch tần suất cao, tạo ra lượng dữ liệu sổ lệnh rất lớn trong thời gian thực. Các hành vi thao túng như spoofing, layering hoặc wash trading có thể xuất hiện trong thời gian rất ngắn, gây sai lệch tín hiệu cung cầu và làm méo mó quyết định của nhà đầu tư.

Dự án này xây dựng một hệ thống giám sát thao túng thị trường theo hướng lai ghép: dùng C++ cho phần lõi sổ lệnh tốc độ cao, Python cho điều phối, Machine Learning cho bộ lọc nhanh, Apache Spark cho phân tán kịch bản rủi ro, và phân tích đồ thị/Shapley Value để truy tìm tài khoản có vai trò trung tâm trong cụm giao dịch nghi vấn.

\section{Mục tiêu}

Các mục tiêu chính của dự án gồm:

\begin{itemize}
    \item Mô phỏng và cập nhật trạng thái Limit Order Book với hiệu năng cao.
    \item Phát hiện nhanh trạng thái sổ lệnh có dấu hiệu spoofing bằng Machine Learning.
    \item Dùng thuật toán tìm kiếm Alpha-Beta và mô phỏng kịch bản để đánh giá rủi ro sâu hơn.
    \item Phân tán hàng trăm kịch bản rủi ro bằng Spark RDD và tổng hợp kết quả bằng Spark SQL.
    \item Phát hiện cụm wash trading bằng NetworkX.
    \item Tính Shapley Value để xác định tài khoản có đóng góp cao nhất, được xem là candidate ringleader cần kiểm tra thêm.
    \item Lưu toàn bộ kết quả thực nghiệm thành file JSON/CSV để phục vụ báo cáo và kiểm chứng.
\end{itemize}

\chapter{Kiến Trúc Hệ Thống}

\section{Mô hình 4 tầng}

\begin{longtable}{p{3cm}p{5cm}p{6cm}}
\toprule
\textbf{Tầng} & \textbf{Thành phần} & \textbf{Vai trò} \\
\midrule
Layer 1 & Machine Learning Fast Filter & Trích xuất đặc trưng từ LOB và dự đoán xác suất spoofing. \\
Layer 2 & C++ LOB + Scenario Search & Mô phỏng sổ lệnh và đánh giá rủi ro bằng heuristic search/Alpha-Beta Search. \\
Layer 3 & Spark RDD + Spark SQL & Phân tán nhiều kịch bản và tổng hợp thống kê rủi ro. \\
Layer 4 & Graph + Shapley Value & Phát hiện wash trading community và xác định candidate ringleader. \\
\bottomrule
\end{longtable}

\section{Luồng xử lý tổng quan}

\begin{lstlisting}
Input / Replay / Redis Stream
        |
        v
Python Orchestrator
        |
        +-- ML Filter
        +-- C++ LOB State
        +-- Spark RDD Scenario Analysis
        +-- Spark SQL Summary
        +-- Graph Community Detection
        +-- Shapley Value Candidate Ringleader Detection
        |
        v
JSON / CSV Output Artifacts
\end{lstlisting}

\chapter{Thiết Kế Lõi C++ Limit Order Book}

\section{Cấu trúc dữ liệu}

Lõi C++ được triển khai trong các file \texttt{src/cpp/lob.h}, \texttt{src/cpp/lob.cpp} và \texttt{src/cpp/bindings.cpp}. Sổ lệnh sử dụng \texttt{std::map} để quản lý hai phía bid và ask. Bid được sắp xếp giảm dần theo giá, còn ask được sắp xếp tăng dần theo giá. Các order được lưu trong \texttt{std::unordered\_map<uint64\_t, Order>} để tra cứu nhanh theo ID.

\section{Chức năng chính}

Lớp \texttt{LOB} hỗ trợ các thao tác:

\begin{itemize}
    \item \texttt{add\_order}: thêm limit order.
    \item \texttt{cancel\_order}: hủy order theo ID.
    \item \texttt{execute\_market\_order}: khớp lệnh market vào phía đối ứng.
    \item \texttt{get\_best\_bid}, \texttt{get\_best\_ask}: lấy giá tốt nhất.
    \item \texttt{get\_spread}: tính spread.
    \item \texttt{get\_cancellation\_rate}: tính tỷ lệ hủy.
    \item \texttt{evaluate\_state}: tính điểm heuristic rủi ro.
    \item \texttt{alpha\_beta\_search}: tìm kiếm/scoring kịch bản rủi ro theo heuristic search.
\end{itemize}

\section{Hàm đánh giá rủi ro}

Hàm heuristic kết hợp spread, imbalance và cancellation rate:

\[
Risk = Spread \times 10 + Imbalance \times 50 + CancellationRate \times 30
\]

Nếu sổ lệnh rơi vào trạng thái crossed book, hệ thống gán điểm rủi ro rất cao để phản ánh bất thường nghiêm trọng.

\section{Đặc tả tìm kiếm Alpha-Beta}

Trong phạm vi project, \texttt{alpha\_beta\_search} được dùng như một lớp heuristic search để đánh giá các trạng thái LOB sinh ra từ kịch bản mô phỏng. Trạng thái \(s\) gồm tập lệnh bid/ask hiện tại, best bid, best ask, spread, imbalance và cancellation rate. Hàm tiện ích cuối cùng là:

\[
U(s)=Risk(s)
\]

với \(Risk(s)\) được tính bởi công thức heuristic ở trên. Các hành động được xét trong mô phỏng gồm thêm lệnh, hủy lệnh và thực hiện market order với volume thay đổi theo từng scenario. Ở tầng Spark, mỗi scenario được xem là một nhánh độc lập; Alpha-Beta Search chỉ dùng để lan truyền điểm heuristic trong cây trạng thái nhỏ của từng scenario, không dùng để chứng minh chiến lược tối ưu tuyệt đối của một tác nhân thao túng.

Mô tả rút gọn:

\begin{lstlisting}
alpha_beta(state, depth, alpha, beta, maximizing):
    if depth == 0 or terminal(state):
        return evaluate_state(state)

    if maximizing:
        value = -inf
        for action in legal_actions(state):
            value = max(value, alpha_beta(apply(state, action),
                                         depth - 1, alpha, beta, false))
            alpha = max(alpha, value)
            if alpha >= beta:
                break
        return value

    value = inf
    for action in legal_actions(state):
        value = min(value, alpha_beta(apply(state, action),
                                     depth - 1, alpha, beta, true))
        beta = min(beta, value)
        if alpha >= beta:
            break
    return value
\end{lstlisting}

\chapter{Machine Learning Fast Filter}

\section{Đặc trưng đầu vào}

Tầng ML trích xuất 5 đặc trưng: spread, bid volume, ask volume, imbalance và cancellation rate. Các đặc trưng được dùng để dự đoán xác suất spoofing bằng Random Forest.

Để đảm bảo khả năng tái lập, các đặc trưng được hiểu theo các quy ước sau. Với \(p^{ask}_1\) là best ask, \(p^{bid}_1\) là best bid, \(q^{bid}_i\) và \(q^{ask}_i\) là khối lượng ở mức giá thứ \(i\):

\[
\mathrm{spread}=p^{ask}_1-p^{bid}_1
\]

\[
V_b=\sum_{i=1}^{L} q^{bid}_i,\qquad
V_a=\sum_{i=1}^{L} q^{ask}_i
\]

\[
\mathrm{imbalance}=\frac{V_b-V_a}{V_b+V_a}
\]

\[
\mathrm{cancel\_rate}=\frac{N_{cancel}}{N_{new}+N_{cancel}}
\]

Trong demo hiện tại, \(V_b\) và \(V_a\) được tính trên các level đang có trong snapshot LOB được đưa vào pipeline. \(\mathrm{imbalance}\) là tỷ lệ có dấu: giá trị dương cho biết phía bid đang chiếm ưu thế về volume, giá trị âm cho biết phía ask chiếm ưu thế. \(\mathrm{cancel\_rate}\) được tính trên cửa sổ sự kiện đang được replay/stream vào trạng thái LOB tại thời điểm đánh giá.

Dữ liệu FI-2010 được chuyển sang replay JSON, sau đó các episode spoofing tổng hợp được chèn vào bằng cách tạo cụm lệnh volume lớn gần best bid/best ask và hủy sau một số event ngắn. Nhãn spoofing được gán cho các event nằm trong episode chèn thêm. Để tránh rò rỉ dữ liệu, train và test dùng file nguồn khác nhau; tập test được tạo từ file raw riêng, không lấy lại các cửa sổ replay gần trùng với tập train.

\section{Model}

Model đã train được lưu tại:

\begin{lstlisting}
models/spoofing_rf_model.pkl
\end{lstlisting}

Metric mẫu:

\begin{longtable}{lr}
\toprule
\textbf{Metric} & \textbf{Giá trị} \\
\midrule
Accuracy & 0.9673 \\
Precision & 0.3867 \\
Recall & 0.5079 \\
F1-score & 0.4391 \\
ROC-AUC & 0.8109 \\
\bottomrule
\end{longtable}

Precision không quá cao vì bài toán có dữ liệu mất cân bằng, nhưng recall và ROC-AUC cho thấy model có khả năng lọc tín hiệu bất thường để chuyển sang tầng phân tích sâu.

\chapter{Spark RDD và Spark SQL}

\section{Phân tán kịch bản}

Khi ML phát hiện rủi ro cao, hệ thống gọi \texttt{analyze\_scenarios\_distributed()} trong \texttt{spark\_engine.py}. Spark RDD tạo nhiều scenario ngẫu nhiên bằng cách thay đổi market order volume, noise orders, giá và volume quanh best bid/best ask.

Mỗi scenario clone trạng thái LOB ban đầu, sau đó gọi C++ heuristic search/Alpha-Beta Search để tính risk score. Vì mỗi scenario là một giả định độc lập về diễn biến thị trường, điểm risk ở tầng này nên được hiểu là điểm heuristic để xếp hạng rủi ro, không phải xác suất đã hiệu chỉnh.

\section{Spark SQL}

Kết quả scenario được chuyển thành Spark DataFrame và đăng ký dưới view \texttt{scenario\_results}. Spark SQL tính tổng số scenario, rủi ro trung bình, rủi ro lớn nhất, rủi ro nhỏ nhất, độ lệch chuẩn và phân vị 95\%.

\chapter{Graph Analytics và Shapley Value}

\section{Phát hiện wash trading}

Các giao dịch tài khoản được biểu diễn thành đồ thị có hướng bằng NetworkX:

\[
G=(V,E,w,t)
\]

trong đó \(V\) là tập tài khoản, \(E\) là giao dịch có hướng từ tài khoản bán sang tài khoản mua, \(w\) là volume giao dịch và \(t\) là thời điểm hoặc thứ tự event. Hệ thống tìm các strongly connected components để phát hiện cụm giao dịch chéo. Kết quả này chỉ cho thấy các tài khoản có quan hệ giao dịch hai chiều hoặc theo chu trình trong cùng cụm; nó chưa tự chứng minh wash trading, ý định thao túng hoặc quyền kiểm soát chung. Vì vậy, cụm được gắn nhãn là nhóm nghi vấn và cần được kiểm tra cùng volume, tần suất, timing và ngữ cảnh giao dịch.

\section{Shapley Value}

Sau khi có community nghi vấn, Spark RDD phân tán nhiều hoán vị Monte Carlo để ước lượng Shapley Value. Hàm đặc trưng \(v(S)\) của một coalition \(S\) được hiểu là tổng đóng góp thanh khoản/giao dịch nội bộ mà nhóm \(S\) tạo ra trong community. Với một tài khoản \(i\), Shapley Value ước lượng đóng góp biên trung bình của \(i\) khi được thêm vào các coalition khác nhau:

\[
\phi_i=\mathbb{E}_{\pi}\left[v(S_{\pi,i}\cup\{i\})-v(S_{\pi,i})\right]
\]

trong đó \(S_{\pi,i}\) là tập tài khoản đứng trước \(i\) trong hoán vị Monte Carlo \(\pi\). Tài khoản có Shapley Value cao nhất được gọi là highest-contribution account hoặc candidate ringleader, không phải kết luận pháp lý cuối cùng về chủ mưu nếu chưa có bằng chứng bổ sung về ý định, quyền kiểm soát hoặc phối hợp ngoài dữ liệu giao dịch.

\chapter{Triển Khai và Môi Trường}

\section{Môi trường khuyến nghị}

\begin{itemize}
    \item Windows với Python 3.11 64-bit.
    \item Java 17.
    \item PySpark 3.5.6.
    \item CMake.
    \item MSVC x64.
    \item Redis qua Docker Desktop.
\end{itemize}

Python 3.12 có thể gây lỗi worker trên PySpark Windows, do đó project dùng Python 3.11 cho demo và streaming.

\section{Build C++}

Module C++ được build thành:

\begin{lstlisting}
build-py311-v2/lob_core.cp311-win_amd64.pyd
\end{lstlisting}

Việc build lại không thay đổi source C++ gốc; nó chỉ tạo artifact để Python import được.

\section{Xử lý vấn đề Windows}

Trên Windows, Spark/Hadoop yêu cầu \texttt{winutils.exe}. Project dùng stub local:

\begin{lstlisting}
hadoop/bin/winutils.exe
\end{lstlisting}

Khi chạy streaming, cần set:

\begin{lstlisting}
HADOOP_HOME=D:\big_data_project\hadoop
\end{lstlisting}

\chapter{Thực Nghiệm}

\section{Demo tĩnh}

Demo tĩnh chạy toàn bộ pipeline và sinh output artifacts tại:

\begin{lstlisting}
outputs/<run_id>/
\end{lstlisting}

Các file quan trọng gồm:

\begin{itemize}
    \item \texttt{01\_initial\_lob\_snapshot.json}
    \item \texttt{02\_ml\_filter\_result.json}
    \item \texttt{03\_scenario\_results.csv}
    \item \texttt{04\_spark\_summary.json}
    \item \texttt{05\_top5\_scenarios.json}
    \item \texttt{06\_wash\_trading\_community.json}
    \item \texttt{07\_shapley\_values.json}
    \item \texttt{10\_final\_result.json}
\end{itemize}

\section{Kết quả mẫu}

Run mẫu:

\begin{lstlisting}
outputs/20260620_083945
\end{lstlisting}

Kết quả ML:

\begin{longtable}{lr}
\toprule
\textbf{Thông tin} & \textbf{Giá trị} \\
\midrule
Spread & 10.0 \\
Bid volume & 50050 \\
Ask volume & 20 \\
Imbalance & 0.9992 \\
Cancellation rate & 0.0 \\
Probability spoofing & 1.0 \\
Inference time & 1.097 ms \\
\bottomrule
\end{longtable}

Kết quả Spark:

\begin{longtable}{lr}
\toprule
\textbf{Metric} & \textbf{Giá trị} \\
\midrule
Total scenarios & 200 \\
Average risk & 149.26 \\
Max risk & 151.67 \\
Min risk & 147.46 \\
Standard deviation & 0.70 \\
P95 risk & 150.12 \\
Elapsed time & 8.1281 seconds \\
\bottomrule
\end{longtable}

Top scenario rủi ro cao nhất:

\begin{longtable}{rrrr}
\toprule
\textbf{Scenario ID} & \textbf{Risk Score} & \textbf{Market Volume} & \textbf{Noise Orders} \\
\midrule
35 & 151.67 & 1017 & 3 \\
183 & 151.61 & 1933 & 2 \\
75 & 151.26 & 991 & 2 \\
16 & 151.00 & 1689 & 2 \\
81 & 150.97 & 607 & 3 \\
\bottomrule
\end{longtable}

Kết quả Shapley:

\begin{longtable}{lr}
\toprule
\textbf{Account} & \textbf{Shapley Value} \\
\midrule
ACC\_RINGLEADER & 453.416667 \\
ACC\_3 & 273.533333 \\
ACC\_1 & 149.966667 \\
ACC\_2 & 123.083333 \\
\bottomrule
\end{longtable}

Kết luận cuối:

\begin{lstlisting}
status: SPOOFING_AND_WASH_TRADING
severity: CRITICAL
ringleader: ACC_RINGLEADER
avg_risk: 149.26
max_risk: 151.67
prob_spoofing: 1.0
\end{lstlisting}

Trường \texttt{ringleader} trong output là tên key của artifact demo; trong diễn giải báo cáo, giá trị này được hiểu là \texttt{candidate ringleader} hoặc tài khoản có đóng góp cao nhất theo Shapley Value. Đây là tín hiệu ưu tiên điều tra, không phải kết luận pháp lý độc lập.

\section{Chính sách severity và xác suất}

Điểm risk của hệ thống là điểm heuristic tổng hợp từ LOB, scenario search và kết quả mô phỏng. Để output có thể kiểm tra lại, severity được ánh xạ theo quy ước:

\begin{longtable}{lr}
\toprule
\textbf{Severity} & \textbf{Điều kiện risk score} \\
\midrule
LOW & \(Risk < 40\) \\
MEDIUM & \(40 \leq Risk < 80\) \\
HIGH & \(80 \leq Risk < 120\) \\
CRITICAL & \(Risk \geq 120\) \\
\bottomrule
\end{longtable}

Trong run mẫu, \texttt{avg\_risk = 149.26} và \texttt{max\_risk = 151.67}, do đó trạng thái được gắn \texttt{CRITICAL}. Trường \texttt{prob\_spoofing} là điểm xác suất do Random Forest trả về từ \texttt{predict\_proba}. Giá trị này chưa được hiệu chỉnh bằng Platt scaling hoặc isotonic regression, vì vậy nó nên được hiểu là score của model để xếp hạng mức nghi vấn, không phải xác suất pháp lý tuyệt đối. Khi cần vận hành thực tế, nên bổ sung calibration và lưu thêm version/checksum của artifact \texttt{10\_final\_result.json}.

\chapter{Đánh Giá}

\section{Ưu điểm}

\begin{itemize}
    \item Kết hợp nhiều tầng phân tích, tránh phụ thuộc vào một tín hiệu duy nhất.
    \item C++ giúp phần mô phỏng LOB và tìm kiếm rủi ro chạy nhanh hơn Python thuần.
    \item Spark cho phép phân tán nhiều kịch bản song song.
    \item Output artifacts giúp chứng minh pipeline đã chạy đầy đủ.
    \item Thiết kế module tách biệt: C++, ML, Spark, Redis, Shapley.
\end{itemize}

\section{Hạn chế}

\begin{itemize}
    \item Streaming trên Windows cần cấu hình thêm Java, Python 3.11 và Hadoop winutils.
    \item Model ML hiện vẫn phụ thuộc nhiều vào dữ liệu tổng hợp hoặc dữ liệu chuyển đổi.
    \item Precision còn thấp, cần cải thiện cân bằng dữ liệu và đặc trưng.
    \item Redis và Spark package làm môi trường chạy phức tạp hơn demo tĩnh.
\end{itemize}

\section{Hướng phát triển}

\begin{itemize}
    \item Tăng dữ liệu thực nghiệm từ FI-2010 hoặc dữ liệu thị trường thật.
    \item Thử các model như XGBoost, Isolation Forest hoặc sequence model.
    \item Tách pipeline streaming thành Docker Compose hoàn chỉnh.
    \item Lưu kết quả streaming thành output artifacts giống demo tĩnh.
    \item Tối ưu C++ core bằng benchmark và profiling.
\end{itemize}

\chapter{Kết Luận}

Dự án đã xây dựng được một hệ thống phát hiện thao túng thị trường có kiến trúc tương đối đầy đủ: từ sổ lệnh C++ hiệu năng cao, bộ lọc ML, mô phỏng rủi ro bằng Spark RDD, tổng hợp Spark SQL, đến phát hiện wash trading và candidate ringleader bằng Shapley Value. Điểm mạnh của hệ thống là khả năng tạo bằng chứng thực nghiệm rõ ràng thông qua các file output JSON/CSV, giúp việc kiểm tra, báo cáo và trình bày kết quả trở nên minh bạch.

Kết quả demo cho thấy pipeline có thể phát hiện tình huống spoofing nghiêm trọng, đánh dấu rủi ro đỏ, nhận diện cụm giao dịch chéo và xác định \texttt{ACC\_RINGLEADER} là tài khoản có đóng góp lớn nhất trong cụm theo Shapley Value. Đây là nền tảng hợp lý để tiếp tục mở rộng thành hệ thống giám sát dữ liệu thị trường thời gian thực.

\end{document}
